
\documentclass[10pt,a4paper]{article}
\usepackage[top=12mm,bottom=12mm,left=14mm,right=14mm,headsep=2mm,footskip=7mm]{geometry}
\usepackage{fontspec}
\usepackage{xeCJK}
\setmainfont{Noto Sans}
\setsansfont{Noto Sans}
\setCJKmainfont{Noto Sans CJK SC}
\setCJKsansfont{Noto Sans CJK SC}
\usepackage{xcolor}
\usepackage{graphicx}
\usepackage{tikz}
\usetikzlibrary{calc,arrows.meta,positioning}
\usepackage{fancyhdr}
\usepackage{hyperref}
\usepackage{ragged2e}
\usepackage{microtype}
\definecolor{P2Paper}{HTML}{FCFAF7}
\definecolor{P2Ink}{HTML}{2D3238}
\definecolor{P2Muted}{HTML}{70757A}
\definecolor{P2BlueLight}{HTML}{D3EAF5}
\definecolor{P2Apricot}{HTML}{F5D9B8}
\definecolor{P2Rose}{HTML}{E8C7D3}
\definecolor{P2Violet}{HTML}{D7D3EA}
\definecolor{P2Blue}{HTML}{5B9FBE}
\definecolor{P2ApricotAccent}{HTML}{D49757}
\definecolor{P2RoseAccent}{HTML}{B97589}
\definecolor{P2VioletAccent}{HTML}{8177A8}
\definecolor{P2Rule}{HTML}{DDE3E7}
\definecolor{P2Panel}{HTML}{FFFDFC}
\pagecolor{P2Paper}
\color{P2Ink}
\setlength{\parindent}{0pt}
\setlength{\parskip}{3.5pt}
\linespread{1.04}
\hypersetup{hidelinks}
\pagestyle{fancy}
\fancyhf{}
\renewcommand{\headrulewidth}{0pt}
\fancyfoot[L]{\fontsize{6.6}{8}\selectfont\color{P2Muted}PROCESS REPORT · WORKFLOW CONSTRUCTION}
\fancyfoot[R]{\fontsize{6.6}{8}\selectfont\color{P2Muted}\thepage}
\newcommand{\eyebrow}[1]{{\fontsize{7.8}{9.4}\selectfont\bfseries\color{P2RoseAccent}#1}\par\vspace{1pt}}
\newcommand{\pagetitle}[1]{{\fontsize{18.5}{21.5}\selectfont\bfseries\color{P2Ink}#1}\par\vspace{3pt}}
\newcommand{\deck}[1]{{\fontsize{9.0}{13.0}\selectfont\color{P2Ink}\RaggedRight #1}\par\vspace{4pt}}
\newcommand{\tracehead}[2]{{\fontsize{7.1}{8.4}\selectfont\bfseries\color{#1}#2}\par}
\newcommand{\tracetitle}[1]{{\fontsize{10.2}{12.2}\selectfont\bfseries\color{P2Ink}#1}\par\vspace{1.5pt}}
\newcommand{\tracebody}[1]{{\fontsize{8.0}{11.2}\selectfont\color{P2Ink}\RaggedRight #1}\par}
\newcommand{\provenance}[1]{\vspace{2pt}{\color{P2Rule}\hrule height .35pt}\vspace{2.5pt}{\fontsize{6.9}{9.2}\selectfont\color{P2Muted}\RaggedRight\textbf{Original Evidence.} #1}\par}
\tikzset{
 userA/.style={rounded corners=1mm,draw=P2ApricotAccent,line width=.55pt,fill=P2Apricot,fill opacity=.34},
 userR/.style={rounded corners=1mm,draw=P2RoseAccent,line width=.55pt,fill=P2Rose,fill opacity=.34},
 gptB/.style={rounded corners=1mm,draw=P2Blue,line width=.55pt,fill=P2BlueLight,fill opacity=.27},
 gptV/.style={rounded corners=1mm,draw=P2VioletAccent,line width=.55pt,fill=P2Violet,fill opacity=.27},
 arrA/.style={-{Stealth[length=2.0mm,width=1.45mm]},draw=P2RoseAccent,line width=.72pt},
 arrB/.style={-{Stealth[length=2.0mm,width=1.45mm]},draw=P2Blue,line width=.72pt},
 arrV/.style={-{Stealth[length=2.0mm,width=1.45mm]},draw=P2VioletAccent,line width=.72pt},
 tagA/.style={circle,minimum size=4.8mm,inner sep=0pt,draw=P2ApricotAccent,line width=.6pt,fill=P2Paper,font=\fontsize{6.7}{7}\selectfont\bfseries},
 tagR/.style={circle,minimum size=4.8mm,inner sep=0pt,draw=P2RoseAccent,line width=.6pt,fill=P2Paper,font=\fontsize{6.7}{7}\selectfont\bfseries},
 tagB/.style={circle,minimum size=4.8mm,inner sep=0pt,draw=P2Blue,line width=.6pt,fill=P2Paper,font=\fontsize{6.7}{7}\selectfont\bfseries}
}
\begin{document}

\eyebrow{WORKFLOW CONSTRUCTION · 14}
\pagetitle{可视化从“装饰”升级为研究表达与证据}
\deck{我确实希望智慧城市作业里的图更丰富，但不是“多画几张高级图”就结束。我的要求是：图必须和真实数据、方法或结论对应，能够帮助我解释异常、参数、方法比较或最终结果。}
\begin{tikzpicture}[x=1mm,y=1mm]
  \node[anchor=north west,inner sep=0] (u) at (0,0) {\includegraphics[width=62mm]{assets/user_visual.png}};
  \node[anchor=north west,inner sep=0] (g) at (0,-42) {\includegraphics[width=70mm]{assets/gpt_visual_core.png}};
  \begin{scope}[shift={(u.south west)},x={($(u.south east)-(u.south west)$)},y={($(u.north west)-(u.south west)$)}]
    \path[userA] (0.42,0.60) rectangle (0.99,0.99); \coordinate (uvis) at (0.99,0.795);
  \end{scope}
  \begin{scope}[shift={(g.south west)},x={($(g.south east)-(g.south west)$)},y={($(g.north west)-(g.south west)$)}]
    \path[gptB] (0.17,0.43) rectangle (0.55,0.57); \coordinate (gcore) at (0.55,0.50);
    \path[gptV] (0.17,0.31) rectangle (0.95,0.38); \coordinate (gthree) at (0.95,0.345);
  \end{scope}
  \draw[arrA] (uvis) .. controls (80,-4) and (80,-65) .. (gcore);
  \draw[arrB] (uvis) .. controls (85,-8) and (85,-90) .. (gthree);
  \node[tagA] at ($(uvis)+(3.8mm,0)$) {1};
  \node[anchor=north west,text width=88mm,align=left] at (91,-2) {
    \tracehead{P2ApricotAccent}{MICRO TRACE 18}\tracetitle{“多一些可视化，但必须与真实数据对应”}\tracebody{我要求增加可视化，同时明确技术图不能靠生成式图片代替。GPT 后来把这件事整理成 visual-first：正式结果优先由真实数据、GIS和绘图库生成。}
    \vspace{7mm}
    \tracehead{P2Blue}{THREE-LAYER SYSTEM}\tracetitle{Publication quality + 时空表达 + 解释层}
    \tracebody{普通科研图负责把数值说清楚；地图、轨迹、OD和时空结构负责体现智慧城市数据的特点；annotation和caption则负责告诉老师这张图到底回答什么问题。}
  };
\end{tikzpicture}
\provenance{上方是我当时对智慧城市可视化提出的真实要求；下方为GPT随后整理出的可视化定位和三层表达方式。}
\newpage
\eyebrow{WORKFLOW CONSTRUCTION · 15}
\pagetitle{Visual Design System：高级不是“复杂”，而是语义统一}
\deck{图一多，我发现如果每次都临时选字体、颜色和布局，最后整份报告会很散。所以我们开始把字体层级、颜色含义、地图底图、annotation和跨图一致性统一起来。}
\begin{tikzpicture}[x=1mm,y=1mm]
  \node[anchor=north west,inner sep=0] (g) at (0,0) {\includegraphics[width=78mm]{assets/gpt_visual_system.png}};
  \begin{scope}[shift={(g.south west)},x={($(g.south east)-(g.south west)$)},y={($(g.north west)-(g.south west)$)}]
    \path[gptB] (0.17,0.63) rectangle (0.95,0.73); \coordinate (gsys) at (0.95,0.68);
    \path[gptV] (0.17,0.25) rectangle (0.55,0.39); \coordinate (gsem) at (0.55,0.32);
  \end{scope}
  \node[anchor=north west,text width=91mm,align=left] at (86,-2) {
    \tracehead{P2Blue}{SYSTEM RULE}\tracetitle{图不再是单张作品，而是同一视觉系统中的证据}
    \tracebody{GPT把我对“专业、统一、可读”的要求整理成一套视觉规则：字号层级尽量固定，颜色有明确含义，地图底图保持克制，annotation只标真正重要的信息。}
    \vspace{8mm}
    \tracehead{P2VioletAccent}{READABILITY}\tracetitle{禁止为了高级感牺牲可读性}
    \tracebody{我后来也把一个原则定得很明确：不要为了“高级感”使用彩虹色、无意义3D或复杂图型堆叠。图看起来丰富没有问题，但必须让老师一眼能看懂。}
  };
\end{tikzpicture}
\provenance{该页来自真实Visual System形成过程；没有把后来锁定的P2配色倒写成我一开始就已经提出的要求。}
\newpage
\eyebrow{WORKFLOW CONSTRUCTION · 16}
\pagetitle{Signature Visualization：每个主要任务至少有一张“解释这个任务”的图}
\deck{最后我没有规定每个任务必须画多少张“高级图”，而是换成一个更实用的要求：每个主要任务至少设计一张最能解释这个任务的 Signature Visualization，而且可读性优先。}
\begin{tikzpicture}[x=1mm,y=1mm]
  \node[anchor=north west,inner sep=0] (g) at (0,0) {\includegraphics[width=78mm]{assets/gpt_signature.png}};
  \begin{scope}[shift={(g.south west)},x={($(g.south east)-(g.south west)$)},y={($(g.north west)-(g.south west)$)}]
    \path[gptB] (0.17,0.845) rectangle (0.97,0.96); \coordinate (gsig) at (0.97,0.90);
    \path[gptV] (0.17,0.34) rectangle (0.74,0.49); \coordinate (gread) at (0.74,0.415);
  \end{scope}
  \node[anchor=north west,text width=91mm,align=left] at (86,-2) {
    \tracehead{P2Blue}{SIGNATURE RULE}\tracetitle{图型围绕研究问题选择，而不是围绕“高级”选择}
    \tracebody{OD flow、Pareto、bivariate map、交互式轨迹图都可以作为候选，但只有真的适合当前问题时才使用。}
    \vspace{8mm}
    \tracehead{P2VioletAccent}{QUALITY GATE}\tracetitle{“看起来丰富，但读起来非常清楚”}
    \tracebody{多维关系可以用 Parallel Coordinates 或 coordinated views；但如果图很“高级”却让老师看不懂，我宁可换回更清楚的表达。}
    \vspace{8mm}
    \tracehead{P2RoseAccent}{REPORT DIFFERENCE}\tracetitle{Experiment与Process的视觉职责也被区分}
    \tracebody{Experiment Report主要让老师看数据、方法和结果；Process Report则更强调真实互动、decision timeline/tree和workflow。两份文档可以长得像同一套设计，但不能用同一套信息结构。}
  };
\end{tikzpicture}
\provenance{该页使用真实讨论中“Signature Visualization、可读性、Experiment/Process视觉分工”的连续内容。}
\end{document}
